\documentclass[output=paper]{langscibook}
\ChapterDOI{10.5281/zenodo.22880238}
\title{\textit{O atovo} -- The sago palm}
\author{Ruth Sima Rigamu and Joyce Maion}
\abstract{}
\begin{document}
\maketitle\label{chap:6}
\begin{figure}[H]
    \captionsetup[subfigure]{oneside,margin={0cm,0cm}}
    \subfloat[]{\includegraphics[trim=0 0 0 0,clip,width=.39\linewidth]{figures/6_1.jpg}\label{fig:6_1}}\hfil
    \subfloat[]{\includegraphics[trim=0 0 20 0,clip,width=.61\linewidth]{figures/6_2.jpg}\label{fig:6_2}}
    \caption{(a) \textit{O atovo} -- The sago palm, (b) \textit{Vua atovo} -- Sago palm nuts}
\end{figure}
\begin{xltabular}{\textwidth}{lXX}
    (\ref{ex:text2-1})&\textit{\textbf{O atovo}}&\textbf{The sago palm}\\
    (\ref{ex:text2-2})&\textit{O naono vai o atovo na nebaneba komana bata ni riori.}&The sago palm is planted.\\
    (\ref{ex:text2-3})&\textit{Ei kou o naono bona eve to vaa'ave riori bona maa inu teori.}&Because it is this palm that they use for thatching their houses.\\
    (\ref{ex:text2-4})&\textit{O atovo mee, be vua, a vua nae repaa kuu, ore paa pura komana vahaa.}&Also, when the sago palm bears nuts, (and) a nut falls down, then it grows by itself again.\\
    (\ref{ex:text2-5})&\textit{Beve sun rova nana tea hatana, na pura vakis nana.}&Even when it is standing in the swamp, it still grows.\\
\end{xltabular}
\begin{figure}[H]
    \captionsetup[subfigure]{oneside,margin={0cm,0cm}}
    \captionsetup{justification=centering}
    \subfloat[]{\includegraphics[trim=20 0 20 0,clip,width=.6\linewidth]{figures/6_3.jpg}\label{fig:6_3}}\hfil
    \subfloat[]{\includegraphics[trim=0 20 0 50,clip,width=.4\linewidth]{figures/6_4.jpg}\label{fig:6_4}}
    \caption{(a) \textit{A kisi atovo} -- A bundle of sago palm leaves,\\ (b) \textit{Huha atovo} – Sewing the thatch}
\end{figure}
\begin{xltabular}{\textwidth}{lXX}
    (\ref{ex:text2-6})&\textit{\textbf{Kosi atovo}}&\textbf{Cutting off the leaves of the sago palm}\\
    (\ref{ex:text2-7})&\textit{Bean tau kosi o atovo, taba momohu toro paku an, ean toro tagava roho o atovo.}&When you are about to cut off the leaves of a sago palm, the first thing you must do, you must fell a sago palm.\\
    (\ref{ex:text2-8})&\textit{Ean repaa mosi hapehape o atovo.}&Then you cut off the branches of the sago palm.\\
    (\ref{ex:text2-9})&\textit{Ean repaa kosi amaa paka atovo kahi bono sinoo vaa tobina.}&Then you cut the sago palm leaves off the stalk in the middle.\footnote{Compare Chapter \ref{chap:5}:\ref{ex:text01-109}.}\\
    (\ref{ex:text2-10})&\textit{Murinae, ean repaa varomo nia maa paka atovo.}&After that, you stack up the sago palm leaves.\\
    (\ref{ex:text2-11})&\textit{Bean varomo vakavara, ean repaa kisi vatataunun riori.}&When you have stacked them up, you tie them together.\\
    (\ref{ex:text2-12})&\textit{Ean repaa taneo mana tea dee voosu amaa kisi atovo vo vaan. 	}&Then you eventually start carrying the bundles of sago palm leaves home to the village.\\
    (\ref{ex:text2-13})&\textit{Bea kavara na maa kisi atovo tei vaan,}&When all the bundles of sago palm leaves are in the village,\\
    (\ref{ex:text2-14})&\textit{ean repaa taneo tea siri sinoo,}&you start tearing off the bones,\\
    (\ref{ex:text2-15})&\textit{o sinoo vaa tobina no paka atovo to siri noman.}&it is the bone in the middle of the sago palm leaf that you tear off (the midrib).\\
    (\ref{ex:text2-16})&\textit{Bean siri vakavara a peha kisi atovo,}&When you have finished tearing off (the midribs of) one bundle of sago palm leaves,\\
    (\ref{ex:text2-17})&\textit{ean repaa vaoromo vaha bana amaa paka atovo.}&you stack up the sago palm leaves again.\\
    (\ref{ex:text2-18})&\textit{Murinae ean repaa kaku eve, a vaamatatopo vaa tea huha.}&After that you fold them in preparation for sewing the thatch.\footnote{See Chapter \ref{chap:17} for details.}\\
\end{xltabular}
\begin{xltabular}{\textwidth}{lXX}
    (\ref{ex:text2-19})&\textit{\textbf{O sikiri}}&\textbf{The midrib}\\
    (\ref{ex:text2-20})&\textit{Amaa sikiri to paa siri vuru an tea maa paka atovo,}&The midribs that you have earlier torn off the sago palm leaves,\\
    (\ref{ex:text2-21})&\textit{eve to pasi huha atovo vahaa ni noman.}&they are what you will use for sewing the sago palm leaves, i.e. the thatch sheets.\\
    (\ref{ex:text2-22})&\textit{Amaa sikiri bona mee, na veveo raara bona maa nou vaa burihina.}&These midribs we also use for sweeping the rubbish in the backyard.\footnote{See Chapter \ref{chap:5}:\ref{ex:text01-110}.}\\
\end{xltabular}
\begin{xltabular}{\textwidth}{lXX}
    (\ref{ex:text2-23})&\textit{\textbf{O rom atovo}}&\textbf{The stem of the sago palm}\\
    (\ref{ex:text2-24})&\textit{O rom atovo na romrom riori, eori repaa kopakopaa bari.}&They cut the stem of the sago palm crosswise into logs, and then they cut them lengthwise.\\
    (\ref{ex:text2-25})&\textit{Murinae eori repaa taraha varaaramun bona komanae bara toon nae,}&After that they shave their inside and their outside to make them smooth,\\
    (\ref{ex:text2-26})&\textit{ahiki ta taba re tahi ribari.}&so that nothing would prick them.\\
    (\ref{ex:text2-27})&\textit{Beve kavara, eori repaa deedee voosu mana bari vo vaan}&When it is finished, they eventually carry them home to the village\\
    (\ref{ex:text2-28})&\textit{mene paku bono baranda na inu.}&for making the floor of the house.\\
\end{xltabular}
\begin{figure}[H]
    \includegraphics[width=.85\linewidth]{figures/6_5.jpg}
    \caption{\textit{A inu gorogoroho} – A sleeping house (under construction)}
    \label{fig:6_5}
\end{figure}
\begin{xltabular}{\textwidth}{lXX}
    (\ref{ex:text2-29})&\textit{\textbf{Vaa'ave na inu}}&\textbf{Thatching the house}\\
    (\ref{ex:text2-30})&\textit{Eori paa huha vakavara vaevuru bona paka atovo vai to paa kosi vuru maori,}&They have already sewed all (the thatch sheets of) sago palm leaves that they had cut off earlier,\\
    (\ref{ex:text2-31})&\textit{matatopo vai tea vaa'ave bona inu.}&in preparation for thatching the house.\\
    (\ref{ex:text2-32})&\textit{Evehee beori mei rori tea vaa'ave toro rahi obai maa roho,}&But before they do the thatching, they must get calamus canes,\footnote{See Chapter \ref{chap:11}:\ref{ex:text7-44}ff.}\\
    (\ref{ex:text2-33})&\textit{bara toro mosi maa roho ta maa naono vai to pasi sun riori beori tau vaa'ave.}&and cut the poles (for the scaffold) on which they will be standing when they are going to thatch the house.\\
    (\ref{ex:text2-34})&\textit{O dorana o naono pasi vaasun vori,}&They will put up four poles,\\
    (\ref{ex:text2-35})&\textit{eori repaa vaobete vananata bono meho dorana beori kisi batari bari bono kara obai.}&and then they lay another four poles across, while tying them to the standing poles with the calamus string.\\
    (\ref{ex:text2-36})&\textit{Eve a poka bona to pasi sun bata nom ori}&This is the scaffold on which they will stand,\\
    (\ref{ex:text2-37})&\textit{beori vaa'ave batari.}&when they are doing the thatching.\\
    (\ref{ex:text2-38})&\textit{Bea poka kavara, taneo vai tea vaa'ave.}&When the scaffold is finished, they start thatching.\\
    (\ref{ex:text2-39})&\textit{A kukan gea dorana aba pasi sunu tea poka,}&Three or four people will stand of the scaffold,\\
    (\ref{ex:text2-40})&\textit{bea meha kukan hee hiava bona vana atovo.}&while another three pass the thatch-sheets up to them.\\
    (\ref{ex:text2-41})&\textit{A dorana aba vai gina to pasi kisi bata ni nom bona vana atovo tea varuvaru na inu.}&The four people on top (of the scaffold) are the ones that will tie the thatch sheets to the rafters of the house.\footnote{See Chapter \ref{chap:16}:\ref{ex:text11-25}ff.}\\
\end{xltabular}
\begin{xltabular}{\textwidth}{lXX}
    (\ref{ex:text2-42})&\textit{\textbf{A huun no atovo o aniani}}&\textbf{The edible sago palm mash}\\
    (\ref{ex:text2-43})&\textit{A abana saka tagava pinopino haari bono atovo vai to pasi koro riori mene ani.}&The people do not carelessly fell the sago palm that they will carve out for eating (the sago).\\
    (\ref{ex:text2-44})&\textit{Na tara rori bono atovo vai o vahuusu vua,}&They see a sago palm close to bearing nuts,\\
    (\ref{ex:text2-45})&\textit{eori repaa tagava bari}&and then they fell it,\\
    (\ref{ex:text2-46})&\textit{ei kou na nata vorien o atovo bona o mataa bara o manoto.}&because they know that this sago palm is good and ripe.\\
    (\ref{ex:text2-47})&\textit{Beori tagava bari, eori repaa kopaa pete bari.}&When they have felled it, they cut it lengthwise.\\
    (\ref{ex:text2-48})&\textit{Murinae eori repaa koro bari bara tasu vamuramuraka bata bari.}&After that they cut it out and mash it (i.e. what they have carved out).\\
    (\ref{ex:text2-49})&\textit{Eori repaa garuhu bari ruene.}&Then they wash it in fresh water.\\
    (\ref{ex:text2-50})&\textit{Beori garuhu bari, ore paa posi batari bari vo tea maa nahu vai amaa beera ge pagete.}&When they have washed it, they squeeze the water out of the mash into large pots or buckets.\\
    (\ref{ex:text2-51})&\textit{O hum vai to ani riori pasi taviana vaavaha komana bata nom.}&The part that they eat gets separated.\\
    (\ref{ex:text2-52})&\textit{Beve raka ore paa kakaavoo.}&When it is dry, it becomes white.\\
\end{xltabular}
\begin{xltabular}{\textwidth}{lXX}
    (\ref{ex:text2-53})&\textit{\textbf{Moroko vagegeana}}&\textbf{Conclusion}\\
    (\ref{ex:text2-54})&\textit{O atovo o mataa, a haihai teve a beera sana en teara to vigu inu ni rara o atovo.}&The sago palm is good; its help is very big for us who use the sago palm for building houses.\\
\end{xltabular}
\end{document}